\chapter{Complex Numbers}

\section{The Algebra of Complex Numbers}

\begin{definition}{Complex Number}{complex_number}
A \textbf{complex number} $z$ is an expression of the form
\begin{equation}
z = a + bi,
\end{equation}
where $a$ and $b$ are real numbers ($a, b \in \mathbb{R}$), and $i$ is the imaginary unit satisfying $i^2 = -1$.
\begin{itemize}
    \item $a = \Re(z)$ is called the \textbf{real part} of $z$.
    \item $b = \Im(z)$ is called the \textbf{imaginary part} of $z$.
\end{itemize}
Thus, we write:
\begin{equation}
z = \Re(z) + i \Im(z).
\end{equation}
If $a = 0$, $z$ is called a \textbf{pure imaginary number}.
\end{definition}

\begin{definition}{Equality of Complex Numbers}{complex_equality}
Two complex numbers $z_1 = a + bi$ and $z_2 = c + di$ are \textbf{equal}, denoted $a + bi = c + di$, if and only if
\begin{equation}
a = c \quad \text{and} \quad b = d.
\end{equation}
Equality of two complex numbers is therefore equivalent to a pair of real equations.
\end{definition}

\begin{proposition}{Basic Algebraic Operations}{alg_operations}
Let $z_1 = a + bi$ and $z_2 = c + di$. The standard arithmetic operations are defined as:
\begin{itemize}
    \item \textbf{Addition and Subtraction}:
    \begin{equation}
    (a + bi) \pm (c + di) := (a \pm c) + (b \pm d)i.
    \end{equation}
    \item \textbf{Multiplication}:
    \begin{equation}
    (a + bi)(c + di) := (ac - bd) + (bc + ad)i,
    \end{equation}
    which follows from the distributive law and $i^2 = -1$.
    \item \textbf{Division}: If $c^2 + d^2 \neq 0$ (i.e., $z_2 \neq 0$), division is defined by rationalizing the denominator:
    \begin{equation}
    \frac{a + bi}{c + di} := \frac{(a + bi)(c - di)}{(c + di)(c - di)} = \left(\frac{ac + bd}{c^2 + d^2}\right) + \left(\frac{bc - ad}{c^2 + d^2}\right)i.
    \end{equation}
\end{itemize}
\end{proposition}

\begin{proposition}{Algebraic Properties of $\mathbb{C}$}{alg_properties}
The set of all complex numbers, denoted by $\mathbb{C}$, forms a field under addition and multiplication. In particular:
\begin{itemize}
    \item \textbf{Commutativity}: $z_1 + z_2 = z_2 + z_1$ and $z_1 z_2 = z_2 z_1$.
    \item \textbf{Associativity}: $(z_1 + z_2) + z_3 = z_1 + (z_2 + z_3)$ and $(z_1 z_2)z_3 = z_1(z_2 z_3)$.
    \item \textbf{Distributivity}: $z_1(z_2 + z_3) = z_1 z_2 + z_1 z_3$.
    \item \textbf{Identities}: $0 = 0 + 0i$ is the additive identity; $1 = 1 + 0i$ is the multiplicative identity.
    \item \textbf{Inverses}: Each $z = a + bi$ has an additive inverse $-z = -a - bi$. Each $z \neq 0$ has a unique multiplicative inverse $z^{-1} = 1/z$.
\end{itemize}
\end{proposition}

\begin{remark}[Lack of Natural Ordering]
Unlike the real numbers $\mathbb{R}$, there is \textbf{no natural ordering} on $\mathbb{C}$ compatible with its field operations. Statements such as $2 + 3i < 3 + 2i$ or $i > 0$ are mathematically meaningless.
\end{remark}

\begin{definition}{Absolute Value (Modulus)}{modulus}
The \textbf{absolute value} or \textbf{modulus} of a complex number $z = a + bi$ is denoted by $|z|$ and defined as:
\begin{equation}
|z| := \sqrt{a^2 + b^2}.
\end{equation}
\begin{itemize}
    \item $|z|$ is always a non-negative real number: $|z| \ge 0$.
    \item $|z| = 0$ if and only if $z = 0$.
\end{itemize}
\end{definition}

\begin{example}{Modulus Calculations}{modulus_ex}
\begin{itemize}
    \item $|0| = 0$,
    \item $\left|\frac{i}{2}\right| = \frac{1}{2}$,
    \item $|3 - 4i| = \sqrt{3^2 + (-4)^2} = \sqrt{9 + 16} = 5$.
\end{itemize}
\end{example}

\begin{proposition}{Distance and Circles in $\mathbb{C}$}{distance_circles}
Let $z_1 = a_1 + b_1 i$ and $z_2 = a_2 + b_2 i$.
\begin{itemize}
    \item The modulus of the difference represents the Euclidean distance between $(a_1, b_1)$ and $(a_2, b_2)$:
    \begin{equation}
    |z_1 - z_2| = |(a_1 - a_2) + (b_1 - b_2)i| = \sqrt{(a_1 - a_2)^2 + (b_1 - b_2)^2}.
    \end{equation}
    \item For a fixed point $z_0 \in \mathbb{C}$ and $r > 0$, the equation of a circle centered at $z_0$ with radius $r$ is:
    \begin{equation}
    |z - z_0| = r.
    \end{equation}
\end{itemize}
\end{proposition}


\section{Complex Conjugates}

\begin{definition}{Complex Conjugate}{conjugate}
The \textbf{complex conjugate} of $z = a + bi$ is denoted by $\bar{z}$ (or $z^*$) and defined by:
\begin{equation}
\bar{z} := a - bi.
\end{equation}
Geometrically, $\bar{z}$ is the reflection of the point $z$ across the real axis.
\end{definition}

\begin{example}{Complex Conjugates}{conj_ex}
\begin{itemize}
    \item $\overline{-1 + 5i} = -1 - 5i$,
    \item $\overline{\pi - i} = \pi + i$,
    \item $\overline{8} = 8$.
\end{itemize}
\end{example}

\begin{proposition}{Properties of the Complex Conjugate}{conj_properties}
For all $z, z_1, z_2 \in \mathbb{C}$:
\begin{enumerate}
    \item \textbf{Real characterization}:
    \begin{equation}
    z = \bar{z} \iff z \in \mathbb{R}.
    \end{equation}
    \item \textbf{Sum and difference}:
    \begin{equation}
    \overline{z_1 \pm z_2} = \bar{z}_1 \pm \bar{z}_2.
    \end{equation}
    \item \textbf{Multiplication}:
    \begin{equation}
    \overline{z_1 z_2} = \bar{z}_1 \bar{z}_2.
    \end{equation}
    \item \textbf{Division}:
    \begin{equation}
    \overline{\left(\frac{z_1}{z_2}\right)} = \frac{\bar{z}_1}{\bar{z}_2} \quad (z_2 \neq 0).
    \end{equation}
    \item \textbf{Real and Imaginary parts}:
    \begin{equation}
    \Re(z) = \frac{z + \bar{z}}{2}, \qquad \Im(z) = \frac{z - \bar{z}}{2i}.
    \end{equation}
    \item \textbf{Involution}:
    \begin{equation}
    \overline{(\bar{z})} = z.
    \end{equation}
    \item \textbf{Modulus}:
    \begin{equation}
    |z| = |\bar{z}|.
    \end{equation}
    \item \textbf{Product with conjugate}:
    \begin{equation}
    z\bar{z} = |z|^2.
    \end{equation}
    \item \textbf{Reciprocal and quotient}:
    \begin{equation}
    \frac{1}{z} = \frac{\bar{z}}{|z|^2} \quad (z \neq 0), \qquad \frac{z_1}{z_2} = \frac{z_1 \bar{z}_2}{|z_2|^2} \quad (z_2 \neq 0).
    \end{equation}
\end{enumerate}
\end{proposition}


\section{Vectors}

\begin{proposition}{Vector Representation and Triangle Inequalities}{vectors_triangle}
A complex number $z = x + iy$ can be identified with the two-dimensional vector $\mathbf{v} = (x, y)$ in the plane, having Euclidean norm $\|\mathbf{v}\| = |z|$.
Addition of complex numbers corresponds to vector addition via the parallelogram law.

For any two complex numbers $z_1, z_2 \in \mathbb{C}$:
\begin{enumerate}
    \item \textbf{Triangle Inequality I}:
    \begin{equation}
    |z_1 + z_2| \le |z_1| + |z_2|.
    \end{equation}
    \item \textbf{Triangle Inequality II (Reverse Triangle Inequality)}:
    \begin{equation}
    |z_2 - z_1| \ge \big||z_2| - |z_1|\big| \ge |z_2| - |z_1|.
    \end{equation}
\end{enumerate}
\end{proposition}


\section{Polar Coordinates}

\begin{definition}{Polar Coordinates and Polar Form}{polar_coords}
A non-zero complex number $z = x + iy$ can be expressed in terms of polar coordinates $(r, \theta)$:
\begin{equation}
x = r \cos\theta, \qquad y = r \sin\theta,
\end{equation}
where:
\begin{itemize}
    \item $r = |z| = \sqrt{x^2 + y^2} > 0$ is the \textbf{modulus}.
    \item $\theta$ is the \textbf{polar angle} (in radians), uniquely determined up to integer multiples of $2\pi$ by:
    \begin{equation}
    \cos\theta = \frac{x}{|z|}, \qquad \sin\theta = \frac{y}{|z|}.
    \end{equation}
\end{itemize}
The \textbf{polar form} of $z$ is:
\begin{equation}
z = r(\cos\theta + i\sin\theta) =: r \cis\theta.
\end{equation}
\end{definition}

\begin{remark}[Quadrant Adjustment for $\arctan$]
While $\tan\theta = y/x$ whenever $x \neq 0$, the relation $\theta = \arctan(y/x)$ only yields the correct angle for points in the 1st and 4th quadrants. For points in the 2nd and 3rd quadrants, one must add or subtract $\pi$:
\begin{equation}
\theta = \begin{cases}
\arctan(y/x) & \text{if } x > 0, \\
\arctan(y/x) + \pi & \text{if } x < 0 \text{ and } y \ge 0, \\
\arctan(y/x) - \pi & \text{if } x < 0 \text{ and } y < 0, \\
+\frac{\pi}{2} & \text{if } x = 0 \text{ and } y > 0, \\
-\frac{\pi}{2} & \text{if } x = 0 \text{ and } y < 0.
\end{cases}
\end{equation}
\end{remark}

\begin{definition}{The Argument Function, Branches, and Principal Value}{argument_branch}
For $z \neq 0$:
\begin{itemize}
    \item The \textbf{argument function} $\arg(z)$ is multi-valued, defined as:
    \begin{equation}
    \arg(z) := \theta + 2k\pi \quad (k \in \mathbb{Z}).
    \end{equation}
    \item A \textbf{branch} of $\arg(z)$ is any single-valued selection taking values in a half-open interval of length $2\pi$, denoted $(\tau, \tau + 2\pi]$:
    \begin{equation}
    \arg_\tau(z) \in (\tau, \tau + 2\pi].
    \end{equation}
    \item The \textbf{principal value} of the argument, denoted $\Arg(z)$, corresponds to the branch taking values in $(-\pi, \pi]$:
    \begin{equation}
    \Arg(z) := \arg_{-\pi}(z) \in (-\pi, \pi].
    \end{equation}
    Consequently:
    \begin{equation}
    \arg(z) = \Arg(z) + 2k\pi \quad (k \in \mathbb{Z}).
    \end{equation}
    \item \textbf{Branch Cut}: $\Arg(z)$ is discontinuous across the negative real axis (a discontinuity jump of $2\pi$). This line of discontinuity is called the \textbf{branch cut}.
    \item The value $\arg(0)$ is not defined for any branch.
\end{itemize}
\end{definition}

\begin{example}{Other Branches of $\arg(z)$}{branch_ex}
\begin{itemize}
    \item $\arg_0(z)$ takes values in $(0, 2\pi]$ and has its branch cut along the positive real axis.
    \item $\arg_{2\pi}(z)$ takes values in $(2\pi, 4\pi]$ with its branch cut along the positive real axis.
\end{itemize}
\end{example}

\begin{proposition}{Geometric Interpretation of Multiplication and Division}{geom_mult_div}
Let $z_1 = r_1 \cis\theta_1$ and $z_2 = r_2 \cis\theta_2$.
\begin{enumerate}
    \item \textbf{Product}:
    \begin{equation}
    z_1 z_2 = r_1 r_2 \big[\cos(\theta_1 + \theta_2) + i \sin(\theta_1 + \theta_2)\big].
    \end{equation}
    Hence:
    \begin{align}
    |z_1 z_2| &= |z_1| |z_2|, \\
    \arg(z_1 z_2) &= \arg(z_1) + \arg(z_2) \pmod{2\pi}.
    \end{align}
    \item \textbf{Quotient}: For $z_2 \neq 0$:
    \begin{equation}
    \frac{z_1}{z_2} = \frac{r_1}{r_2} \big[\cos(\theta_1 - \theta_2) + i \sin(\theta_1 - \theta_2)\big].
    \end{equation}
    Hence:
    \begin{align}
    \left|\frac{z_1}{z_2}\right| &= \frac{|z_1|}{|z_2|}, \\
    \arg\left(\frac{z_1}{z_2}\right) &= \arg(z_1) - \arg(z_2) \pmod{2\pi}.
    \end{align}
    \item \textbf{Conjugate and Reciprocal}:
    \begin{equation}
    \arg(\bar{z}) = -\arg(z) \pmod{2\pi}, \qquad \arg\left(\frac{1}{z}\right) = -\arg(z) \pmod{2\pi}.
    \end{equation}
    The vectors $\bar{z}$ and $z^{-1}$ point in the same direction (parallel vectors).
\end{enumerate}
\end{proposition}


\section{The Complex Exponential}

\begin{definition}{The Complex Exponential and Euler's Formula}{complex_exp}
For a purely imaginary number $iy$, \textbf{Euler's formula} states:
\begin{equation}
e^{iy} := \cos y + i \sin y.
\end{equation}
For a general complex number $z = x + iy$, the \textbf{complex exponential} is defined by:
\begin{equation}
e^z := e^{x+iy} = e^x e^{iy} = e^x (\cos y + i \sin y).
\end{equation}
\end{definition}

\begin{proposition}{Properties of the Complex Exponential}{exp_properties}
\begin{enumerate}
    \item \textbf{Multiplicative Property}:
    \begin{equation}
    e^{z_1} e^{z_2} = e^{z_1 + z_2}.
    \end{equation}
    \item \textbf{Division Rule}:
    \begin{equation}
    \frac{e^{z_1}}{e^{z_2}} = e^{z_1 - z_2}.
    \end{equation}
    \item \textbf{Derivative}:
    \begin{equation}
    \frac{d}{dz} e^z = e^z.
    \end{equation}
    \item \textbf{Modulus and Argument}:
    \begin{equation}
    |e^z| = e^{\Re(z)} = e^x, \qquad \arg(e^z) = \Im(z) + 2k\pi = y + 2k\pi.
    \end{equation}
    In particular, for any real $\theta$:
    \begin{equation}
    |e^{i\theta}| = 1.
    \end{equation}
    \item \textbf{Periodicity}: $e^z$ is periodic with purely imaginary period $2\pi i$:
    \begin{equation}
    e^{z + 2\pi i} = e^z \quad \forall z \in \mathbb{C}.
    \end{equation}
\end{enumerate}
\end{proposition}

\begin{proposition}{Exponential Form and Key Identities}{exp_identities}
Using Euler's formula, any non-zero complex number $z$ can be written in \textbf{exponential polar form}:
\begin{equation}
z = r e^{i\theta} = |z| e^{i \arg(z)}.
\end{equation}
\begin{itemize}
    \item \textbf{Key Values}:
    \begin{equation}
    e^{i0} = 1, \quad e^{i\pi/2} = i, \quad e^{-i\pi/2} = -i, \quad e^{i\pi} = -1, \quad e^{2k\pi i} = 1 \quad (k \in \mathbb{Z}).
    \end{equation}
    \item \textbf{Trigonometric Relations}:
    \begin{equation}
    \cos\theta = \Re(e^{i\theta}) = \frac{e^{i\theta} + e^{-i\theta}}{2}, \qquad \sin\theta = \Im(e^{i\theta}) = \frac{e^{i\theta} - e^{-i\theta}}{2i}.
    \end{equation}
    \item \textbf{Multiplication and Division}:
    \begin{equation}
    z_1 z_2 = (r_1 e^{i\theta_1})(r_2 e^{i\theta_2}) = (r_1 r_2) e^{i(\theta_1 + \theta_2)}, \qquad \frac{z_1}{z_2} = \frac{r_1 e^{i\theta_1}}{r_2 e^{i\theta_2}} = \left(\frac{r_1}{r_2}\right) e^{i(\theta_1 - \theta_2)}.
    \end{equation}
    \item \textbf{Conjugation}:
    \begin{equation}
    \bar{z} = \overline{r e^{i\theta}} = r e^{-i\theta}.
    \end{equation}
\end{itemize}
\end{proposition}

\begin{theorem}{De Moivre's Formula}{de_moivre}
For any integer $n \in \mathbb{Z}$ and real number $\theta \in \mathbb{R}$:
\begin{equation}
(e^{i\theta})^n = e^{i n\theta},
\end{equation}
which is equivalent to:
\begin{equation}
(\cos\theta + i \sin\theta)^n = \cos(n\theta) + i \sin(n\theta).
\end{equation}
\end{theorem}


\section{Powers and Roots}

\subsection{Powers}

\begin{proposition}{Integer Powers of a Complex Number}{integer_powers}
Let $z = r e^{i\theta} = r(\cos\theta + i\sin\theta)$. For any integer $n \in \mathbb{Z}$:
\begin{equation}
z^n = r^n e^{i n\theta} = r^n \big(\cos(n\theta) + i \sin(n\theta)\big).
\end{equation}
In particular:
\begin{align}
z^2 &= r^2 e^{i 2\theta}, \\
z^3 &= r^3 e^{i 3\theta}.
\end{align}
\end{proposition}

\subsection{Roots}

\begin{definition}{$m$-th Roots of a Complex Number}{mth_roots}
Let $z \in \mathbb{C}$ with $z \neq 0$, and let $m \in \mathbb{N}^+$. A complex number $\zeta$ is an \textbf{$m$-th root} of $z$ if:
\begin{equation}
\zeta^m = z.
\end{equation}
Every non-zero complex number $z = r e^{i\theta}$ has \textbf{exactly $m$ distinct $m$-th roots}, denoted $z^{1/m}$, given by:
\begin{equation}
z^{1/m} = \sqrt[m]{r} \, e^{i(\theta + 2k\pi)/m} = \sqrt[m]{|z|} \left[\cos\left(\frac{\theta + 2k\pi}{m}\right) + i \sin\left(\frac{\theta + 2k\pi}{m}\right)\right]
\end{equation}
for $k = 0, 1, 2, \dots, m-1$, where $\sqrt[m]{r}$ denotes the unique positive real $m$-th root of $r$.
\end{definition}

\begin{example}{Roots of $\zeta^4 = 1$}{roots_unity_4}
The four solutions of $\zeta^4 = 1$ are obtained from $1 = e^{i 2k\pi}$ ($k=0,1,2,3$):
\begin{itemize}
    \item $k = 0 \implies \zeta_0 = e^{i 0} = 1$,
    \item $k = 1 \implies \zeta_1 = e^{i\pi/2} = i$,
    \item $k = 2 \implies \zeta_2 = e^{i\pi} = -1$,
    \item $k = 3 \implies \zeta_3 = e^{i 3\pi/2} = -i$.
\end{itemize}
\end{example}

\begin{definition}{Roots of Unity and Primitive Roots}{roots_of_unity}
The $m$ distinct $m$-th roots of unity, denoted $1^{1/m}$, are:
\begin{equation}
1^{1/m} = e^{i 2k\pi/m} = \cos\left(\frac{2k\pi}{m}\right) + i \sin\left(\frac{2k\pi}{m}\right), \quad k = 0, 1, \dots, m-1.
\end{equation}
Defining the principal generator
\begin{equation}
\omega_m := e^{i 2\pi/m} = \cos\left(\frac{2\pi}{m}\right) + i \sin\left(\frac{2\pi}{m}\right),
\end{equation}
the complete set of $m$-th roots of unity can be expressed as:
\begin{equation}
\{1, \omega_m, \omega_m^2, \dots, \omega_m^{m-1}\}.
\end{equation}
\end{definition}

\begin{proposition}{Geometric Interpretation and Sum of Roots}{geom_roots}
\begin{enumerate}
    \item \textbf{Geometric Distribution}: The $m$ roots of unity lie equally spaced along the unit circle, separated by angular increments of $2\pi/m$ radians, forming the vertices of a regular $m$-sided polygon inscribed in the unit circle.
    \begin{itemize}
        \item For $m = 3$ (Cube roots of unity): Vertices of an equilateral triangle:
        \begin{equation}
        1, \quad \omega_3 = e^{i 2\pi/3} = -\frac{1}{2} + i\frac{\sqrt{3}}{2}, \quad \omega_3^2 = e^{i 4\pi/3} = -\frac{1}{2} - i\frac{\sqrt{3}}{2}.
        \end{equation}
        \item For $m = 4$: Vertices of a square: $\{1, i, -1, -i\}$.
        \item For $m = 5$: Vertices of a regular pentagon.
    \end{itemize}
    \item \textbf{Sum of Roots of Unity}: The center of mass of the roots is at the origin:
    \begin{equation}
    1 + \omega_m + \omega_m^2 + \dots + \omega_m^{m-1} = \sum_{k=0}^{m-1} \omega_m^k = 0.
    \end{equation}
\end{enumerate}
\end{proposition}


\section{Planar Sets and Topology in \texorpdfstring{$\mathbb{C}$}{C}}

\begin{definition}{Neighborhoods and Disks}{neighborhoods}
Let $z_0 \in \mathbb{C}$ and $\rho > 0$.
\begin{itemize}
    \item The \textbf{open disk} (or \textbf{circular neighborhood}) of radius $\rho$ centered at $z_0$ is:
    \begin{equation}
    B(z_0, \rho) := \{z \in \mathbb{C} : |z - z_0| < \rho\}.
    \end{equation}
    The \textbf{open unit disk} is the set $\{z \in \mathbb{C} : |z| < 1\}$.
    \item The \textbf{closed disk} of radius $\rho$ centered at $z_0$ is:
    \begin{equation}
    \overline{B}(z_0, \rho) := \{z \in \mathbb{C} : |z - z_0| \le \rho\}.
    \end{equation}
\end{itemize}
\end{definition}

\begin{definition}{Interior Points and Open Sets}{open_sets}
Let $S \subset \mathbb{C}$.
\begin{itemize}
    \item A point $z_0 \in S$ is an \textbf{interior point} of $S$ if there exists some $\rho > 0$ such that the neighborhood $B(z_0, \rho)$ is completely contained within $S$:
    \begin{equation}
    B(z_0, \rho) \subset S.
    \end{equation}
    \item A set $S$ is called an \textbf{open set} if every point of $S$ is an interior point of $S$.
\end{itemize}
\end{definition}

\begin{example}{Open and Non-Open Sets}{open_ex}
\begin{itemize}
    \item Every open disk $B(z_0, \rho)$ is an open set.
    \item Annuli $\rho_1 < |z - z_0| < \rho_2$ are open sets.
    \item Half-planes such as $\Im(z) > 0$ or $\Re(z) > 0$, and strips $1 < \Re(z) < 2$, are open sets.
    \item The set $\{z : |z - 3| \ge 2\}$ is \textbf{not} open, because no point on the circle $|z - 3| = 2$ is an interior point.
    \item An open interval of the real line (e.g., $(0, 1) \subset \mathbb{C}$) is \textbf{not} an open set in $\mathbb{C}$, because it contains no 2D open disk.
\end{itemize}
\end{example}

\begin{definition}{Polygonal Paths and Connectedness (Domains)}{domains}
\begin{itemize}
    \item A \textbf{polygonal path} joining $w_1$ to $w_{n+1}$ is a continuous chain of line segments $l_1, l_2, \dots, l_n$, where each $l_k$ is the line segment joining $w_k$ to $w_{k+1}$.
    \item An open set $S$ is \textbf{connected} if any two points $z_1, z_2 \in S$ can be joined by a polygonal path that lies entirely within $S$.
    \item A non-empty, connected open set is called a \textbf{domain}.
\end{itemize}
\end{definition}

\begin{example}{Connected vs. Disconnected Sets}{conn_ex}
The set $\mathbb{C} \setminus \{z : |z| = 1\}$ (all points in the plane not on the unit circle) is an open set, but it is \textbf{not connected} (and thus not a domain), because any path from a point inside $|z|<1$ to a point outside $|z|>1$ must cross the unit circle.
\end{example}

\begin{theorem}{Zero Partial Derivatives in a Domain}{const_derivative}
Let $u(x, y)$ be a real-valued function with continuous partial derivatives defined on a domain $D$. If
\begin{equation}
\frac{\partial u}{\partial x} = 0 \quad \text{and} \quad \frac{\partial u}{\partial y} = 0 \quad \forall (x, y) \in D,
\end{equation}
then $u(x, y)$ is identically \textbf{constant} throughout $D$:
\begin{equation}
u(x, y) \equiv \text{constant} \quad \text{in } D.
\end{equation}
\end{theorem}

\begin{definition}{Boundary Points and Closed Sets}{boundary_closed}
Let $S \subset \mathbb{C}$.
\begin{itemize}
    \item A point $z_0$ is a \textbf{boundary point} of $S$ if every neighborhood $B(z_0, \rho)$ contains at least one point in $S$ and at least one point not in $S$.
    \item The set of all boundary points of $S$ is called the \textbf{boundary} (or \textbf{frontier}) of $S$, denoted $\partial S$.
    \item A domain contains none of its boundary points ($D \cap \partial D = \emptyset$).
    \item A set $S$ is \textbf{closed} if it contains all of its boundary points ($\partial S \subset S$).
\end{itemize}
\end{definition}

\begin{example}{Boundary and Non-Closed Sets}{bound_ex}
\begin{itemize}
    \item The boundary of the open disk $B(z_0, \rho)$ is the circle $|z - z_0| = \rho$.
    \item The puncture disk $0 < |z| \le 1$ is \textbf{neither} open (it contains boundary points $|z|=1$) \textbf{nor} closed (it does not contain the boundary point $z = 0$).
    \item The closed disk $|z - z_0| \le \rho$ contains its boundary and is closed.
\end{itemize}
\end{example}

\begin{definition}{Bounded, Compact Sets and Regions}{compact_region}
\begin{itemize}
    \item A set $S$ is \textbf{bounded} if there exists $R > 0$ such that $|z| < R$ for all $z \in S$ (i.e., $S$ is contained in some ball around the origin). Otherwise, $S$ is \textbf{unbounded}.
    \item A set $S$ is \textbf{compact} if it is both \textbf{closed} and \textbf{bounded}.
    \item A \textbf{region} is a domain together with some, none, or all of its boundary points. A closed region is the closure of a domain ($\overline{D} = D \cup \partial D$).
\end{itemize}
\end{definition}
